% 2010 年江西卷理科第 12 题：水中的正五角星与四个导函数图象选项。
% 来源 PDF：普通高考/2010/2010江西理.pdf，SHA-256 994a56c33da151e22af3fabe292d1b5b193866dfc35773a4431e6dfc8a0cc934。
% 原页：第 1 页（A3 三栏）第 2 栏中部，书页码 607。
% 重绘依据：照原页，正五角星（角尖向上）浸在圆柱水柱内，星顶在筒口以下；
%   水面为筒内的虚线椭圆，星的两条下臂尖没入水面以下；正确答案 D：S' 正比于
%   水面截线长度，由 0 增到下方凹角处的局部极大、略降、再升到最宽处、最后降到 0，且连续。
% 说明：四幅导函数图为按题设定性特征绘制的示意（原页每幅宽不足 3cm）。
\begin{tikzpicture}[semithick]
  % ---- 上半：圆柱与五角星 ----
  \begin{scope}[shift={(0,0)}]
    % 圆柱：上沿实线椭圆，两侧母线，下沿虚线椭圆
    \draw (0,0) ellipse (1.9 and 0.5);
    \draw (-1.9,0) -- (-1.9,-2.5);
    \draw (1.9,0) -- (1.9,-2.5);
    \draw[dashed] (0,-2.5) ellipse (1.9 and 0.5);
    % 正五角星：外接圆半径 1.12，圆心 (0,-1.2)
    \begin{scope}[shift={(0,-1.2)},scale=1.12]
      \fill[gray!55] (90:1) -- (162:0.382) -- (234:1) -- (306:0.382) -- (18:1)
        -- (54:0.382) -- (126:1) -- (198:0.382) -- (270:1) -- (342:0.382) -- cycle;
      \draw (90:1) -- (162:0.382) -- (234:1) -- (306:0.382) -- (18:1)
        -- (54:0.382) -- (126:1) -- (198:0.382) -- (270:1) -- (342:0.382) -- cycle;
    \end{scope}
    % 水面：筒内虚线椭圆（压在星上，表示没入水中的部分）
    \draw[dashed] (0,-1.55) ellipse (1.9 and 0.5);
  \end{scope}
  % ---- 下半：四个选项 ----
  % (A) 有断点与尖峰
  \begin{scope}[shift={(-0.2,-5.0)}]
    \draw[-stealth] (-0.35,0) -- (2.0,0) node[right] {$t$};
    \draw[-stealth] (0,-0.55) -- (0,1.05) node[above] {$y$};
    \node[below left] at (0,0) {$O$};
    \draw (0,0) -- (0.55,0.28);
    \draw[fill=white] (0.55,0.28) circle (1.4pt);
    \draw[fill=white] (0.62,0.72) circle (1.4pt);
    \draw (0.62,0.72) -- (1.05,0.42) -- (1.28,0.55) -- (1.55,0.1);
    \draw (1.28,0.55) -- (1.44,0.78) -- (1.62,0.1);
    \node at (0.85,-0.95) {（A）};
  \end{scope}
  % (B) 先升后骤降并穿越横轴
  \begin{scope}[shift={(2.9,-5.0)}]
    \draw[-stealth] (-0.35,0) -- (2.0,0) node[right] {$t$};
    \draw[-stealth] (0,-0.85) -- (0,0.8) node[above] {$y$};
    \node[below left] at (0,0) {$O$};
    \draw (0.0,-0.02) -- (0.5,0.22) -- (0.66,0.3);
    \draw (0.66,0.66) -- (0.9,0.34) -- (1.2,-0.62);
    \draw (0.66,0.3) -- (0.66,0.66);
    \draw (1.2,-0.62) -- (1.42,0.1) -- (1.62,-0.12) -- (1.9,0.2);
    \node at (0.85,-1.2) {（B）};
  \end{scope}
  % (C) 单调递增
  \begin{scope}[shift={(6.0,-5.0)}]
    \draw[-stealth] (-0.35,0) -- (2.0,0) node[right] {$t$};
    \draw[-stealth] (0,-0.5) -- (0,1.15) node[above] {$y$};
    \node[below left] at (0,0) {$O$};
    \draw (0,0) -- (0.5,0.3) -- (0.75,0.45) -- (0.78,0.78) -- (1.05,0.9) -- (1.35,1.02);
    \node at (0.85,-0.95) {（C）};
  \end{scope}
  % (D) 升—略降—再升到最高—降到 0
  \begin{scope}[shift={(9.1,-5.0)}]
    \draw[-stealth] (-0.35,0) -- (2.1,0) node[right] {$t$};
    \draw[-stealth] (0,-0.5) -- (0,1.15) node[above] {$y$};
    \node[below left] at (0,0) {$O$};
    \draw (0,0) -- (0.45,0.42) -- (0.62,0.4) -- (1.05,0.95) -- (1.25,0.68) -- (1.55,0.62) -- (1.82,0);
    \node at (0.9,-0.95) {（D）};
  \end{scope}
\end{tikzpicture}
